\section{Weil divisors on a regular curve}
\label{mk-chap:RiemannRoch_WeilDivisor}

Let \(X\) be a Noetherian integral separated scheme that is regular in
codimension one.  Thus every codimension-one local ring of \(X\) is a
discrete valuation ring.  This chapter constructs the group of Weil
divisors, the divisor of a rational function, and linear equivalence.  It
also records the coefficient-degree formalism used for curves over an
algebraically closed field.

The final degree-zero theorem for principal divisors requires the
finite-morphism geometry of complete nonsingular curves.  Its two geometric
inputs are stated separately below.  The later Riemann--Roch argument is developed in
\ref{mk-chap:RiemannRoch_Adelic}.
The divisor conventions follow Hartshorne \cite{hartshorne-ag}.

\section{Prime divisors and Weil divisors}

\begin{definition}[Prime divisor]
  \label{mk-def:prime_divisor}
  \lean{AlgebraicGeometry.Scheme.PrimeDivisor}
  \leanok
  \source{hartshorne-algebraic-geometry:page-0147}
  \dcref{custom}

  A prime divisor of \(X\) is a codimension-one integral closed subscheme.
  Equivalently, it is determined by a point \(x\in X\) whose closure has
  codimension one.  We encode the latter condition by
  \[
    \operatorname{coheight}(x)=1.
  \]
  The reduced closed subscheme on \(\overline{\{x\}}\) is integral, so this
  point and its codimension witness determine the prime divisor.
\end{definition}

\begin{definition}[Weil divisor group]
  \label{mk-def:codim1_cycles}
  \lean{AlgebraicGeometry.Scheme.WeilDivisor}
  \leanok
  \source{hartshorne-algebraic-geometry:page-0147}
  \dcref{custom}

  \uses{mk-def:prime_divisor}

  The group of Weil divisors of \(X\) is the free abelian group on its prime
  divisors:
  \[
    \operatorname{Div}(X)
      =\bigoplus_{Y\in X^{(1)}}\mathbb Z\cdot Y.
  \]
  Thus a divisor is a finite formal sum \(D=\sum_Y n_Y Y\).  It is
  \emph{effective} when every \(n_Y\geq 0\).
\end{definition}

\begin{definition}[Regularity in codimension one]
  \label{mk-def:isRegularInCodimensionOne}
  \lean{AlgebraicGeometry.Scheme.IsRegularInCodimensionOne}
  \leanok
  \source{hartshorne-algebraic-geometry:page-0147}
  \dcref{custom}

  \uses{mk-def:prime_divisor}

  An integral scheme is \emph{regular in codimension one} if
  \(\mathcal O_{X,Y}\) is a discrete valuation ring for every prime divisor
  \(Y\in X^{(1)}\).
\end{definition}

\begin{lemma}[The stalk at a prime divisor is a DVR]
  \label{mk-def:isRegularInCodimOne_dvrStalk}
  \lean{AlgebraicGeometry.Scheme.IsRegularInCodimensionOne.instIsDiscreteValuationRingStalk}
  \leanok
  \dcref{custom}

  \uses{mk-def:isRegularInCodimensionOne, mk-def:prime_divisor}

  If \(X\) is regular in codimension one and \(Y\in X^{(1)}\), then
  \(\mathcal O_{X,Y}\) is a discrete valuation ring.
\end{lemma}
\begin{proof}
  \leanok
  This is the defining property of regularity in codimension one.
\end{proof}

\begin{lemma}[Dimension of a prime-divisor stalk]
  \label{mk-def:isRegularInCodimOne_krullDimStalk}
  \lean{AlgebraicGeometry.Scheme.IsRegularInCodimensionOne.instKrullDimLEStalk}
  \leanok
  \dcref{custom}

  \uses{mk-def:isRegularInCodimOne_dvrStalk}

  If \(X\) is regular in codimension one and \(Y\in X^{(1)}\), then
  \(\dim\mathcal O_{X,Y}\leq 1\).
\end{lemma}
\begin{proof}
  \leanok
  A discrete valuation ring is a principal ideal domain of Krull dimension
  at most one.
\end{proof}

\section{Restriction to an open subscheme}
\label{mk-sec:open_immersion_descent}

Let \(U\subseteq X\) be open.

\begin{lemma}[Extensionality of prime divisors]
  \label{mk-lem:primeDivisor_ext}
  \lean{AlgebraicGeometry.Scheme.PrimeDivisor.ext}
  \leanok
  \dcref{custom}

  \uses{mk-def:prime_divisor}

  Two prime divisors represented by the same point of \(X\) are equal.
\end{lemma}
\begin{proof}
  \leanok
  Their codimension witnesses are propositions, hence proof-irrelevant.
\end{proof}

\begin{definition}[Restriction of a prime divisor]
  \label{mk-def:primeDivisor_restrictToOpen}
  \lean{AlgebraicGeometry.Scheme.PrimeDivisor.restrictToOpen}
  \leanok
  \dcref{custom}

  \uses{mk-def:prime_divisor}

  If \(Y\in X^{(1)}\) and its generic point lies in \(U\), restriction gives
  a prime divisor \(Y|_U\in U^{(1)}\).
\end{definition}

\begin{definition}[Extension from an open subscheme]
  \label{mk-def:primeDivisor_ofOpen}
  \lean{AlgebraicGeometry.Scheme.PrimeDivisor.ofOpen}
  \leanok
  \dcref{custom}

  \uses{mk-def:prime_divisor}

  A prime divisor \(Z\in U^{(1)}\) determines the prime divisor of \(X\)
  represented by the image of its generic point.
\end{definition}

\begin{lemma}[Prime divisors on an open subscheme]
  \label{mk-lem:primeDivisor_equivOpen}
  \lean{AlgebraicGeometry.Scheme.PrimeDivisor.equivOpen}
  \leanok
  \dcref{custom}

  \uses{mk-def:primeDivisor_restrictToOpen, mk-def:primeDivisor_ofOpen}

  Restriction and extension are inverse bijections
  \[
    \{Y\in X^{(1)}\mid \eta_Y\in U\}\;\simeq\;U^{(1)}.
  \]
\end{lemma}
\begin{proof}
  \leanok
  Both composites preserve the representing point.  The result follows from
  \ref{mk-lem:primeDivisor_ext}.
\end{proof}

\begin{lemma}[Point of a restricted prime divisor]
  \label{mk-lem:primeDivisor_restrictToOpen_point}
  \lean{AlgebraicGeometry.Scheme.PrimeDivisor.restrictToOpen_point}
  \leanok
  \dcref{custom}

  \uses{mk-def:primeDivisor_restrictToOpen}

  The point representing \(Y|_U\) is \(\eta_Y\), regarded as a point of \(U\).
\end{lemma}
\begin{proof}
  \leanok
  This follows directly from the definition of restriction.
\end{proof}

\begin{lemma}[Point of an extended prime divisor]
  \label{mk-lem:primeDivisor_ofOpen_point}
  \lean{AlgebraicGeometry.Scheme.PrimeDivisor.ofOpen_point}
  \leanok
  \dcref{custom}

  \uses{mk-def:primeDivisor_ofOpen}

  The point representing the extension of \(Z\in U^{(1)}\) is the image of
  \(\eta_Z\) in \(X\).
\end{lemma}
\begin{proof}
  \leanok
  This follows directly from the definition of extension.
\end{proof}

\begin{lemma}[Stalks under restriction to an open]
  \label{mk-lem:primeDivisor_stalkIso}
  \lean{AlgebraicGeometry.Scheme.PrimeDivisor.stalkIso}
  \leanok
  \dcref{custom}

  \uses{mk-def:primeDivisor_restrictToOpen}

  If \(\eta_Y\in U\), restriction induces a canonical isomorphism
  \[
    \mathcal O_{U,Y|_U}\;\cong\;\mathcal O_{X,Y}.
  \]
\end{lemma}
\begin{proof}
  \leanok
  Stalks of the structure sheaf are unchanged after restriction to an open
  neighbourhood of the point.
\end{proof}

\begin{theorem}[Regularity in codimension one is open-local]
  \label{mk-thm:isRegularInCodimensionOne_open}
  \lean{AlgebraicGeometry.Scheme.IsRegularInCodimensionOne.instOpen}
  \leanok
  \dcref{custom}

  \uses{mk-def:isRegularInCodimensionOne}

  If \(X\) is regular in codimension one, then every integral open subscheme
  \(U\subseteq X\) is regular in codimension one.
\end{theorem}
\begin{proof}
  \leanok
  \uses{mk-def:primeDivisor_ofOpen, mk-lem:primeDivisor_stalkIso}
  Extend a prime divisor of \(U\) to \(X\).  Its stalk in \(X\) is a DVR, and
  \ref{mk-lem:primeDivisor_stalkIso} transports this property to the
  corresponding stalk in \(U\).
\end{proof}

\section{Orders of rational functions}

\subsection{Naturality under isomorphisms}
\label{mk-subsec:ordFrac_ringEquiv}

\begin{lemma}[Order length is invariant under a ring isomorphism]
  \label{mk-lem:ord_ringEquiv}
  \lean{Ring.ord_ringEquiv}
  \leanok
  \dcref{custom}

  Let \(e\colon R\cong S\) be an isomorphism of commutative rings.  For
  \(x\in R\), the order lengths of \(x\) and \(e(x)\) agree.
\end{lemma}
\begin{proof}
  \leanok
  The isomorphism \(e\) induces an equivalence
  \(R/(x)\cong S/(e(x))\), and module length is invariant under this
  equivalence.
\end{proof}

\begin{lemma}[Non-zero-divisors are invariant under isomorphism]
  \label{mk-lem:nonZeroDivisors_ringEquiv}
  \lean{Ring.nonZeroDivisors_ringEquiv}
  \leanok
  \dcref{custom}

  For a ring isomorphism \(e\colon R\cong S\), an element \(r\in R\) is a
  non-zero-divisor if and only if \(e(r)\) is a non-zero-divisor.
\end{lemma}
\begin{proof}
  \leanok
  Multiplication by \(e(r)\) is conjugate, through \(e\), to multiplication
  by \(r\).
\end{proof}

\begin{lemma}[Naturality of the order monoid homomorphism]
  \label{mk-lem:ordMonoidWithZeroHom_ringEquiv}
  \lean{Ring.ordMonoidWithZeroHom_ringEquiv}
  \leanok
  \dcref{custom}

  \uses{mk-lem:ord_ringEquiv, mk-lem:nonZeroDivisors_ringEquiv}

  The order monoid-with-zero homomorphism is invariant under a ring
  isomorphism.
\end{lemma}
\begin{proof}
  \leanok
  Split according to whether the element is a non-zero-divisor.  The
  non-zero-divisor branch is \ref{mk-lem:ord_ringEquiv}; the other branch is
  sent to zero on both sides.
\end{proof}

\begin{lemma}[Naturality of fraction-field order]
  \label{mk-lem:ordFrac_ringEquiv}
  \lean{Ring.ordFrac_ringEquiv}
  \leanok
  \dcref{custom}

  \uses{mk-lem:ordMonoidWithZeroHom_ringEquiv}

  Let \(R\cong S\) be an isomorphism of Noetherian integral domains of
  dimension at most one, and let \(K_R\cong K_S\) be a compatible
  isomorphism of fraction fields.  Then
  \[
    \operatorname{ord}_{S}(e_K(x))
      =\operatorname{ord}_{R}(x)
    \qquad (x\in K_R).
  \]
\end{lemma}
\begin{proof}
  \leanok
  The assertion is immediate for \(x=0\).  Otherwise write \(x=a/b\) with
  \(a\) and \(b\) non-zero-divisors.  Compatibility identifies
  \(e_K(x)\) with \(e(a)/e(b)\), and the result follows by applying
  \ref{mk-lem:ordMonoidWithZeroHom_ringEquiv} to numerator and denominator.
\end{proof}

\begin{definition}[Function field of a nonempty open subscheme]
  \label{mk-def:functionFieldIso}
  \lean{AlgebraicGeometry.Scheme.Opens.functionFieldIso}
  \leanok
  \dcref{custom}

  If \(X\) is integral and \(U\subseteq X\) is nonempty and open, restriction
  induces a canonical isomorphism
  \[
    K(U)\;\cong\;K(X).
  \]
\end{definition}

\begin{lemma}[Compatibility of the function-field isomorphism]
  \label{mk-lem:functionFieldIso_compat}
  \lean{AlgebraicGeometry.Scheme.PrimeDivisor.functionFieldIso_compat}
  \leanok
  \dcref{custom}

  \uses{mk-def:functionFieldIso, mk-lem:primeDivisor_stalkIso}

  Let \(Y\in X^{(1)}\) have generic point in a nonempty integral open
  \(U\).  The diagram
  \[
  \begin{tikzcd}
    \mathcal O_{U,Y|_U} \arrow[r] \arrow[d, "\sim"']
      & K(U) \arrow[d, "\sim"] \\
    \mathcal O_{X,Y} \arrow[r] & K(X)
  \end{tikzcd}
  \]
  commutes.
\end{lemma}
\begin{proof}
  \leanok
  Both composites send a germ represented on an open neighbourhood of
  \(\eta_Y\) to the same germ at the generic point.
\end{proof}

\begin{lemma}[Fraction-field order across an open immersion]
  \label{mk-lem:ordFrac_stalkIso_naturality}
  \lean{AlgebraicGeometry.Scheme.PrimeDivisor.ordFrac_stalkIso_naturality}
  \leanok
  \dcref{custom}

  \uses{mk-lem:primeDivisor_stalkIso, mk-lem:ordFrac_ringEquiv}

  Suppose the local rings in the preceding square are Noetherian domains of
  dimension at most one.  Any compatible isomorphism \(K(U)\cong K(X)\)
  preserves their fraction-field orders.
\end{lemma}
\begin{proof}
  \leanok
  Apply \ref{mk-lem:ordFrac_ringEquiv} to the stalk isomorphism and the
  compatible function-field isomorphism.
\end{proof}

\begin{definition}[Order along a prime divisor]
  \label{mk-def:order_at_point}
  \lean{AlgebraicGeometry.Scheme.RationalMap.order}
  \leanok
  \source{hartshorne-algebraic-geometry:page-0147,
    hartshorne-algebraic-geometry:page-0148}
  \dcref{custom}

  \uses{mk-def:prime_divisor}

  Let \(X\) be integral and locally Noetherian, and let \(Y\in X^{(1)}\)
  have local ring of dimension at most one.  The \emph{order of vanishing}
  of \(f\in K(X)\) along \(Y\) is the integer
  \[
    \operatorname{ord}_Y(f)=v_Y(f).
  \]
  We use the harmless total convention \(\operatorname{ord}_Y(0)=0\).
\end{definition}

\begin{lemma}[Order is invariant under restriction to an open]
  \label{mk-lem:order_eq_order_restrict}
  \lean{AlgebraicGeometry.Scheme.PrimeDivisor.order_eq_order_restrict}
  \leanok
  \dcref{custom}

  \uses{mk-def:order_at_point, mk-def:functionFieldIso,
    mk-def:primeDivisor_restrictToOpen}

  Let \(X\) be integral, locally Noetherian, and regular in codimension one.
  If \(Y\in X^{(1)}\) meets a nonempty integral open \(U\), then for every
  \(f\in K(U)\),
  \[
    \operatorname{ord}_Y(f)
      =\operatorname{ord}_{Y|_U}(f),
  \]
  where the left side uses \(K(U)\cong K(X)\).
\end{lemma}
\begin{proof}
  \leanok
  \uses{mk-lem:functionFieldIso_compat,
    mk-lem:ordFrac_stalkIso_naturality}
  The compatibility square of \ref{mk-lem:functionFieldIso_compat} supplies
  the hypothesis of \ref{mk-lem:ordFrac_stalkIso_naturality}.
\end{proof}

\subsection{Valuation identities}
\label{mk-subsec:order_identities}

\begin{lemma}[Order of zero]
  \label{mk-lem:order_zero}
  \lean{AlgebraicGeometry.Scheme.RationalMap.order_zero}
  \leanok
  \dcref{custom}

  \uses{mk-def:order_at_point}
  For every prime divisor \(Y\), \(\operatorname{ord}_Y(0)=0\).
\end{lemma}
\begin{proof}
  \leanok
  This is the total convention in \ref{mk-def:order_at_point}.
\end{proof}

\begin{lemma}[Order of one]
  \label{mk-lem:order_one}
  \lean{AlgebraicGeometry.Scheme.RationalMap.order_one}
  \leanok
  \dcref{custom}

  \uses{mk-def:order_at_point}
  For every prime divisor \(Y\), \(\operatorname{ord}_Y(1)=0\).
\end{lemma}
\begin{proof}
  \leanok
  The valuation of a multiplicative unit is zero.
\end{proof}

\begin{lemma}[Order of a product]
  \label{mk-lem:order_mul_of_ne_zero}
  \lean{AlgebraicGeometry.Scheme.RationalMap.order_mul_of_ne_zero}
  \leanok
  \dcref{custom}

  \uses{mk-def:order_at_point}
  If \(f,g\in K(X)\) are nonzero, then
  \[
    \operatorname{ord}_Y(fg)
      =\operatorname{ord}_Y(f)+\operatorname{ord}_Y(g).
  \]
\end{lemma}
\begin{proof}
  \leanok
  This is multiplicativity of the discrete valuation, written additively.
\end{proof}

\begin{lemma}[Order of an inverse]
  \label{mk-lem:order_inv}
  \lean{AlgebraicGeometry.Scheme.RationalMap.order_inv}
  \leanok
  \dcref{custom}

  \uses{mk-def:order_at_point}
  For \(f\in K(X)\),
  \[
    \operatorname{ord}_Y(f^{-1})=-\operatorname{ord}_Y(f).
  \]
\end{lemma}
\begin{proof}
  \leanok
  For \(f\neq0\), additivity applied to \(ff^{-1}=1\) gives the result.  The
  total convention makes the equality true also for \(f=0\).
\end{proof}

\begin{lemma}[Order of a unit inverse]
  \label{mk-lem:order_units_inv}
  \lean{AlgebraicGeometry.Scheme.RationalMap.order_units_inv}
  \leanok
  \dcref{custom}

  \uses{mk-lem:order_inv}
  For \(u\in K(X)^\times\),
  \(\operatorname{ord}_Y(u^{-1})=-\operatorname{ord}_Y(u)\).
\end{lemma}
\begin{proof}
  \leanok
  This is \ref{mk-lem:order_inv} applied to the underlying field element.
\end{proof}

\begin{lemma}[Order is insensitive to sign]
  \label{mk-lem:order_neg}
  \lean{AlgebraicGeometry.Scheme.RationalMap.order_neg}
  \leanok
  \dcref{custom}

  \uses{mk-def:order_at_point}
  For \(f\in K(X)\),
  \(\operatorname{ord}_Y(-f)=\operatorname{ord}_Y(f)\).
\end{lemma}
\begin{proof}
  \leanok
  Since \((-f)^2=f^2\), additivity gives twice each side; cancellation in
  \(\mathbb Z\) proves the equality.
\end{proof}

\begin{lemma}[Order of a power]
  \label{mk-lem:order_pow_of_ne_zero}
  \lean{AlgebraicGeometry.Scheme.RationalMap.order_pow_of_ne_zero}
  \leanok
  \dcref{custom}

  \uses{mk-lem:order_one, mk-lem:order_mul_of_ne_zero}
  If \(f\neq0\) and \(n\in\mathbb N\), then
  \[
    \operatorname{ord}_Y(f^n)=n\,\operatorname{ord}_Y(f).
  \]
\end{lemma}
\begin{proof}
  \leanok
  Induct on \(n\), using \ref{mk-lem:order_one} at zero and
  \ref{mk-lem:order_mul_of_ne_zero} at the successor step.
\end{proof}

\section{Closed points and coefficient degree}

\begin{definition}[Divisor associated with a closed point]
  \label{mk-def:divisor_closed_point}
  \lean{AlgebraicGeometry.Scheme.WeilDivisor.ofClosedPoint}
  \leanok
  \dcref{custom}

  \uses{mk-def:codim1_cycles}

  Let \(P\) be a closed point of a scheme \(C\).  Define
  \[
    [P]=
    \begin{cases}
      1\cdot\langle P,h\rangle,&
        \text{if }h:\operatorname{coheight}(P)=1,\\
      0,&\text{otherwise.}
    \end{cases}
  \]
\end{definition}

\begin{lemma}[Closed points of a nonsingular curve are prime divisors]
  \label{mk-lem:curve_closed_point_prime_divisor}
  \dcref{custom}
  \uses{mk-def:divisor_closed_point}

  Let \(C\) be a nonsingular integral curve over an algebraically closed
  field.  Every closed point \(P\in C\) has coheight one, and the divisor
  \([P]\) of \ref{mk-def:divisor_closed_point} is the corresponding prime
  divisor with multiplicity one.
\end{lemma}
\begin{proof}
  A closed point of an integral curve is distinct from the generic point.
  Since the curve has dimension one, the specialization chain from the
  closed point to the generic point has length one.  Thus
  \(\operatorname{coheight}(P)=1\), and the positive branch of
  \ref{mk-def:divisor_closed_point} gives the asserted divisor.
\end{proof}

\begin{lemma}[Closed point in the codimension-one branch]
  \label{mk-lem:ofClosedPoint_eq_single}
  \lean{AlgebraicGeometry.Scheme.WeilDivisor.ofClosedPoint_eq_single}
  \leanok
  \dcref{custom}

  \uses{mk-def:divisor_closed_point}

  If \(h:\operatorname{coheight}(P)=1\), then
  \([P]=1\cdot\langle P,h\rangle\).
\end{lemma}
\begin{proof}
  \leanok
  Select the positive branch in \ref{mk-def:divisor_closed_point}.
\end{proof}

\begin{lemma}[Closed point outside the codimension-one branch]
  \label{mk-lem:ofClosedPoint_eq_zero}
  \lean{AlgebraicGeometry.Scheme.WeilDivisor.ofClosedPoint_eq_zero}
  \leanok
  \dcref{custom}

  \uses{mk-def:divisor_closed_point}

  If \(\operatorname{coheight}(P)\neq1\), then \([P]=0\).
\end{lemma}
\begin{proof}
  \leanok
  Select the negative branch in \ref{mk-def:divisor_closed_point}.
\end{proof}

\begin{definition}[Coefficient degree]
  \label{mk-def:divisor_degree}
  \lean{AlgebraicGeometry.Scheme.WeilDivisor.degree}
  \leanok
  \dcref{custom}

  \uses{mk-def:codim1_cycles}

  For any scheme \(X\), define the coefficient degree of
  \(D=\sum_Y n_YY\) by
  \[
    \deg(D)=\sum_Y n_Y.
  \]
  If \(X\) is a nonsingular curve over an algebraically closed field, this
  is the usual degree because every closed point has residue degree one.
\end{definition}

\begin{theorem}[Coefficient degree is additive]
  \label{mk-thm:divisor_degree_hom}
  \lean{AlgebraicGeometry.Scheme.WeilDivisor.degree_hom}
  \leanok
  \dcref{custom}

  \uses{mk-def:divisor_degree}

  Coefficient degree defines a homomorphism
  \[
    \deg\colon\operatorname{Div}(X)\longrightarrow\mathbb Z.
  \]
\end{theorem}
\begin{proof}
  \leanok
  The sum of coefficients is additive under addition of finite formal sums.
\end{proof}

\begin{lemma}[Bundled degree agrees with coefficient degree]
  \label{mk-lem:degree_hom_apply}
  \lean{AlgebraicGeometry.Scheme.WeilDivisor.degree_hom_apply}
  \leanok
  \dcref{custom}

  \uses{mk-thm:divisor_degree_hom, mk-def:divisor_degree}
  For every divisor \(D\), the homomorphism of
  \ref{mk-thm:divisor_degree_hom} takes \(D\) to \(\deg(D)\).
\end{lemma}
\begin{proof}
  \leanok
  Both sides are the sum of the coefficients of \(D\).
\end{proof}

\begin{lemma}[Degree of zero]
  \label{mk-lem:degree_zero}
  \lean{AlgebraicGeometry.Scheme.WeilDivisor.degree_zero}
  \leanok
  \dcref{custom}

  \uses{mk-def:divisor_degree}
  One has \(\deg(0)=0\).
\end{lemma}
\begin{proof}
  \leanok
  The empty finite sum is zero.
\end{proof}

\begin{lemma}[Degree of a sum]
  \label{mk-lem:degree_add}
  \lean{AlgebraicGeometry.Scheme.WeilDivisor.degree_add}
  \leanok
  \dcref{custom}

  \uses{mk-def:divisor_degree}
  For divisors \(D_1,D_2\),
  \(\deg(D_1+D_2)=\deg(D_1)+\deg(D_2)\).
\end{lemma}
\begin{proof}
  \leanok
  \uses{mk-thm:divisor_degree_hom}
  This is additivity in \ref{mk-thm:divisor_degree_hom}.
\end{proof}

\begin{lemma}[Degree of a negation]
  \label{mk-lem:degree_neg}
  \lean{AlgebraicGeometry.Scheme.WeilDivisor.degree_neg}
  \leanok
  \dcref{custom}

  \uses{mk-def:divisor_degree}
  For a divisor \(D\), \(\deg(-D)=-\deg(D)\).
\end{lemma}
\begin{proof}
  \leanok
  \uses{mk-thm:divisor_degree_hom}
  Apply the homomorphism of \ref{mk-thm:divisor_degree_hom} to \(-D\).
\end{proof}

\begin{lemma}[Degree of a difference]
  \label{mk-lem:degree_sub}
  \lean{AlgebraicGeometry.Scheme.WeilDivisor.degree_sub}
  \leanok
  \dcref{custom}

  \uses{mk-def:divisor_degree}
  For divisors \(D_1,D_2\),
  \(\deg(D_1-D_2)=\deg(D_1)-\deg(D_2)\).
\end{lemma}
\begin{proof}
  \leanok
  \uses{mk-lem:degree_add, mk-lem:degree_neg}
  Combine \ref{mk-lem:degree_add} and \ref{mk-lem:degree_neg}.
\end{proof}

\begin{lemma}[Degree of a one-point divisor]
  \label{mk-lem:degree_single}
  \lean{AlgebraicGeometry.Scheme.WeilDivisor.degree_single}
  \leanok
  \dcref{custom}

  \uses{mk-def:divisor_degree}
  For \(Y\in X^{(1)}\) and \(n\in\mathbb Z\),
  \(\deg(nY)=n\).
\end{lemma}
\begin{proof}
  \leanok
  The divisor \(nY\) has the single coefficient \(n\).
\end{proof}

\section{Principal divisors}

\begin{lemma}[Finite support of the order function]
  \label{mk-lem:rationalMap_order_finite_support}
  \lean{AlgebraicGeometry.rationalMap_order_finite_support}
  \leanok
  \source{hartshorne-algebraic-geometry:page-0148}
  \dcref{custom}

  \uses{mk-def:order_at_point}

  Let \(X\) be a Noetherian integral scheme regular in codimension one.
  For every \(f\in K(X)\), the set
  \[
    \{Y\in X^{(1)}\mid\operatorname{ord}_Y(f)\neq0\}
  \]
  is finite.
\end{lemma}
\begin{proof}
  \leanok
  If \(f=0\), the support is empty by convention.  Otherwise choose a finite
  affine open cover \(X=\bigcup_i\operatorname{Spec}R_i\).  On one chart,
  write \(f=a/b\).  If the height-one prime \(\mathfrak p_Y\) contains
  neither \(a\) nor \(b\), both become units in
  \((R_i)_{\mathfrak p_Y}\), so \(\operatorname{ord}_Y(f)=0\).
  Hence every point of the support is a height-one prime containing \(a\)
  or \(b\).  Such a prime is minimal over \((a)\) or \((b)\), and a
  Noetherian ring has only finitely many minimal primes over an ideal.
  Taking the finite union over the cover proves the claim.
\end{proof}

\begin{remark}
  \label{mk-rem:finite_support_needs_noetherian}
  \dcref{custom}
  \uses{mk-lem:rationalMap_order_finite_support}
  Global Noetherianness cannot be replaced by local Noetherianness.  The
  affine line with infinitely many origins is integral, locally Noetherian,
  and regular in codimension one, but its coordinate function has order one
  at every origin.  The scheme is not quasi-compact.
\end{remark}

\begin{definition}[Principal divisor]
  \label{mk-def:principal_divisor}
  \lean{AlgebraicGeometry.Scheme.WeilDivisor.principal}
  \leanok
  \source{hartshorne-algebraic-geometry:page-0148}
  \dcref{custom}

  \uses{mk-def:codim1_cycles, mk-def:order_at_point,
    mk-lem:rationalMap_order_finite_support}

  For a nonzero \(f\in K(X)\), its principal divisor is
  \[
    \operatorname{div}(f)
      =\sum_{Y\in X^{(1)}}\operatorname{ord}_Y(f)\,Y.
  \]
  The sum is finite by \ref{mk-lem:rationalMap_order_finite_support}.
\end{definition}

\begin{lemma}[Coefficient of a principal divisor]
  \label{mk-lem:principal_apply}
  \lean{AlgebraicGeometry.Scheme.WeilDivisor.principal_apply}
  \leanok
  \dcref{custom}

  \uses{mk-def:principal_divisor}
  For \(Y\in X^{(1)}\),
  \[
    \operatorname{div}(f)(Y)=\operatorname{ord}_Y(f).
  \]
\end{lemma}
\begin{proof}
  \leanok
  This is the coefficient formula in the defining finite sum.
\end{proof}

\begin{lemma}[Principal divisor of one]
  \label{mk-lem:principal_one}
  \lean{AlgebraicGeometry.Scheme.WeilDivisor.principal_one}
  \leanok
  \dcref{custom}

  \uses{mk-def:principal_divisor}
  One has \(\operatorname{div}(1)=0\).
\end{lemma}
\begin{proof}
  \leanok
  \uses{mk-lem:principal_apply, mk-lem:order_one}
  Every coefficient is \(\operatorname{ord}_Y(1)=0\).
\end{proof}

\begin{theorem}[The principal-divisor homomorphism]
  \label{mk-thm:principal_hom}
  \lean{AlgebraicGeometry.Scheme.WeilDivisor.principal_hom}
  \leanok
  \source{hartshorne-algebraic-geometry:page-0148}
  \dcref{custom}

  \uses{mk-def:principal_divisor}

  The assignment
  \[
    K(X)^\times\longrightarrow\operatorname{Div}(X),
    \qquad f\longmapsto\operatorname{div}(f)
  \]
  is a homomorphism from the multiplicative group to the additive group.
\end{theorem}
\begin{proof}
  \leanok
  \uses{mk-lem:principal_apply, mk-lem:order_mul_of_ne_zero, mk-lem:order_one}
  Compare coefficients at every \(Y\).  The product formula for the order
  gives
  \(\operatorname{div}(fg)(Y)=\operatorname{div}(f)(Y)
    +\operatorname{div}(g)(Y)\), and \ref{mk-lem:order_one} treats the
  identity.
\end{proof}

\section{Degree under finite morphisms}

\begin{definition}[Pullback of a divisor along a finite morphism of curves]
  \label{mk-def:finite_curve_divisor_pullback}
  \dcref{custom}
  \uses{mk-lem:curve_closed_point_prime_divisor, mk-def:order_at_point}

  Let \(\varphi\colon C\to D\) be a finite morphism of nonsingular curves
  over an algebraically closed field.
  For a closed point \(Q\in D\), choose a local parameter \(t\) at \(Q\) and
  set
  \[
    \varphi^*[Q]
      =\sum_{\varphi(P)=Q}\operatorname{ord}_P(t)\,[P].
  \]
  This is independent of the local parameter and extends linearly to
  \(\varphi^*\colon\operatorname{Div}(D)\to\operatorname{Div}(C)\).
\end{definition}

\begin{theorem}[Degree of a finite pullback]
  \label{mk-thm:finite_curve_degree_pullback}
  \dcref{custom}
  \uses{mk-def:finite_curve_divisor_pullback, mk-def:divisor_degree}

  If \(\varphi\colon C\to D\) is a finite morphism of complete nonsingular
  curves over an algebraically closed field and
  \(E\in\operatorname{Div}(D)\), then
  \[
    \deg(\varphi^*E)=\deg(\varphi)\deg(E).
  \]
\end{theorem}
\begin{proof}
  It suffices to take \(E=[Q]\).  On an affine neighbourhood
  \(\operatorname{Spec}B\) of \(Q\), let \(A\) be the integral closure of
  \(B\) in \(K(C)\), localize at \(Q\), and write the resulting finite
  torsion-free \(\mathcal O_{D,Q}\)-module as a free module of rank
  \([K(C):K(D)]=\deg(\varphi)\).  If \(t\) is a local parameter at \(Q\),
  the dimension of \(A/tA\) is this rank.  Decomposing \(A/tA\) at the
  points \(P\) over \(Q\) identifies the same dimension with
  \(\sum_{\varphi(P)=Q}\operatorname{ord}_P(t)\), which is
  \(\deg(\varphi^*[Q])\).
\end{proof}

\begin{lemma}[A rational function gives a finite map to the projective line]
  \label{mk-lem:rational_function_finite_map_to_p1}
  \dcref{custom}
  \uses{mk-def:finite_curve_divisor_pullback, mk-def:principal_divisor}

  Let \(C\) be a complete nonsingular curve over an algebraically closed field
  \(k\), and let
  \(f\in K(C)\setminus k\).  There is a finite morphism
  \(\varphi_f\colon C\to\mathbb P^1_k\) inducing the inclusion
  \(k(f)\subseteq K(C)\), and
  \[
    \operatorname{div}(f)
      =\varphi_f^*([0]-[\infty]).
  \]
\end{lemma}
\begin{proof}
  The rational function defines a morphism to \(\mathbb P^1\) because \(C\)
  is complete and nonsingular.  It is nonconstant, hence dominant; the
  induced extension \(K(C)/k(f)\) is finite, so the morphism is finite.
  At every closed point, the pullback multiplicity of \(0\) or \(\infty\)
  is respectively the order of the zero or the pole of \(f\).  Equality of
  all coefficients gives the divisor formula.
\end{proof}

\begin{theorem}[Principal divisors have degree zero]
  \label{mk-thm:principal_deg_zero}
  \lean{AlgebraicGeometry.Scheme.WeilDivisor.principal_degree_zero}
  \source{hartshorne-algebraic-geometry:page-0155}
  \dcref{custom}

  \uses{mk-def:principal_divisor, mk-def:divisor_degree}

  Let \(C\) be a complete nonsingular curve over an algebraically closed
  field.  For every nonzero \(f\in K(C)\),
  \[
    \deg(\operatorname{div}(f))=0.
  \]
\end{theorem}
\begin{proof}
  \uses{mk-lem:rational_function_finite_map_to_p1,
    mk-thm:finite_curve_degree_pullback}
  If \(f\) is constant, its order at every point is zero.  Otherwise
  \ref{mk-lem:rational_function_finite_map_to_p1} gives a finite morphism
  \(\varphi_f\colon C\to\mathbb P^1\) with
  \(\operatorname{div}(f)=\varphi_f^*([0]-[\infty])\).  Therefore
  \[
    \deg(\operatorname{div}(f))
      =\deg(\varphi_f)\deg([0]-[\infty])
      =\deg(\varphi_f)(1-1)=0
  \]
  by \ref{mk-thm:finite_curve_degree_pullback}.
\end{proof}

\section{Positive parts}

\begin{definition}[Positive part of a Weil divisor]
  \label{mk-def:WeilDivisor_positivePart}
  \lean{AlgebraicGeometry.Scheme.WeilDivisor.positivePart}
  \leanok
  \dcref{custom}

  \uses{mk-def:codim1_cycles}

  For \(D=\sum_Y n_YY\), define
  \[
    D_+=\sum_Y\max(n_Y,0)\,Y.
  \]
\end{definition}

\begin{lemma}[Positive part of zero]
  \label{mk-lem:positivePart_zero}
  \lean{AlgebraicGeometry.Scheme.WeilDivisor.positivePart_zero}
  \leanok
  \dcref{custom}

  \uses{mk-def:WeilDivisor_positivePart}
  One has \(0_+=0\).
\end{lemma}
\begin{proof}
  \leanok
  Every coefficient is \(\max(0,0)=0\).
\end{proof}

\begin{lemma}[Positive part of a one-point divisor]
  \label{mk-lem:positivePart_single}
  \lean{AlgebraicGeometry.Scheme.WeilDivisor.positivePart_single}
  \leanok
  \dcref{custom}

  \uses{mk-def:WeilDivisor_positivePart}
  For \(Y\in X^{(1)}\) and \(n\in\mathbb Z\),
  \[
    (nY)_+=\max(n,0)Y.
  \]
\end{lemma}
\begin{proof}
  \leanok
  This follows coefficientwise from the definition.
\end{proof}

\begin{lemma}[Degree of a positive part]
  \label{mk-lem:degree_positivePart_eq_sum_max}
  \lean{AlgebraicGeometry.Scheme.WeilDivisor.degree_positivePart_eq_sum_max}
  \leanok
  \dcref{custom}

  \uses{mk-def:WeilDivisor_positivePart, mk-def:divisor_degree}
  If \(D=\sum_Yn_YY\), then
  \[
    \deg(D_+)=\sum_Y\max(n_Y,0).
  \]
\end{lemma}
\begin{proof}
  \leanok
  Expand the two definitions and compare the finite sums.
\end{proof}

\begin{lemma}[Sum over positive coefficients]
  \label{mk-lem:finsupp_sum_max_zero_eq_sum_filter_pos}
  \lean{Finsupp.sum_max_zero_eq_sum_filter_pos}
  \leanok
  \dcref{custom}

  For a finitely supported function \(a\mapsto n_a\in\mathbb Z\),
  \[
    \sum_a\max(n_a,0)=\sum_{a:\,n_a>0}n_a.
  \]
\end{lemma}
\begin{proof}
  \leanok
  Split each coefficient according to whether it is positive.  Positive
  coefficients contribute themselves to both sides and all other
  coefficients contribute zero.
\end{proof}

\begin{lemma}[A simple zero contributes positive degree]
  \label{mk-lem:one_le_degree_positivePart_principal_of_order_one}
  \lean{AlgebraicGeometry.Scheme.WeilDivisor.one_le_degree_positivePart_principal_of_order_one}
  \leanok
  \dcref{custom}

  \uses{mk-def:principal_divisor, mk-def:WeilDivisor_positivePart,
    mk-def:divisor_degree}

  If \(f\neq0\) and \(\operatorname{ord}_{Y_0}(f)=1\), then
  \[
    1\leq\deg\bigl(\operatorname{div}(f)_+\bigr).
  \]
\end{lemma}
\begin{proof}
  \leanok
  \uses{mk-lem:principal_apply, mk-lem:degree_positivePart_eq_sum_max}
  In the sum of nonnegative coefficients, the \(Y_0\)-term is
  \(\max(\operatorname{ord}_{Y_0}(f),0)=1\); every other term is
  nonnegative.
\end{proof}

\section{Linear equivalence and the divisor class group}

\begin{definition}[Linear equivalence]
  \label{mk-def:linear_equivalence}
  \lean{AlgebraicGeometry.Scheme.WeilDivisor.LinearEquivalence}
  \leanok
  \source{hartshorne-algebraic-geometry:page-0148}
  \dcref{custom}

  \uses{mk-def:principal_divisor}

  Two divisors \(D,D'\in\operatorname{Div}(X)\) are linearly equivalent,
  written \(D\sim D'\), if there is a nonzero \(f\in K(X)\) such that
  \[
    D-D'=\operatorname{div}(f).
  \]
\end{definition}

\begin{lemma}[Linear equivalence is an equivalence relation]
  \label{mk-lem:linear_equivalence_equivalence}
  \dcref{custom}
  \uses{mk-def:linear_equivalence}
  Linear equivalence is reflexive, symmetric, and transitive.
\end{lemma}
\begin{proof}
  \uses{mk-thm:principal_hom, mk-lem:principal_one, mk-lem:order_units_inv}
  Reflexivity uses \(\operatorname{div}(1)=0\).  If
  \(D-D'=\operatorname{div}(f)\), then
  \(D'-D=\operatorname{div}(f^{-1})\), proving symmetry.  Adding the
  equations for \(D\sim D'\) and \(D'\sim D''\), and using
  \(\operatorname{div}(fg)=\operatorname{div}(f)+\operatorname{div}(g)\),
  proves transitivity.
\end{proof}

\begin{definition}[Divisor class group]
  \label{mk-def:divisor_class_group}
  \dcref{custom}
  \uses{mk-thm:principal_hom}
  The divisor class group is
  \[
    \operatorname{Cl}(X)
      =\operatorname{Div}(X)/
        \operatorname{im}\bigl(K(X)^\times
          \xrightarrow{\operatorname{div}}\operatorname{Div}(X)\bigr).
  \]
\end{definition}

\begin{theorem}[Degree on the divisor class group of a curve]
  \label{mk-thm:divisor_class_group_degree}
  \dcref{custom}
  \uses{mk-def:divisor_class_group, mk-def:divisor_degree}
  On a complete nonsingular curve over an algebraically closed field,
  coefficient degree descends to a homomorphism
  \[
    \deg\colon\operatorname{Cl}(C)\longrightarrow\mathbb Z.
  \]
\end{theorem}
\begin{proof}
  \uses{mk-thm:principal_deg_zero}
  By \ref{mk-thm:principal_deg_zero}, the degree homomorphism vanishes on the
  subgroup of principal divisors, hence factors uniquely through the
  quotient.
\end{proof}
